   \documentstyle[%


righttag%


,11pt%


,amssymb%


,russian%               remove in Eng.


,izvru%                 Set izven% in Eng.


% ,izvru-ks             for brief communications


]{amsart}





%       Math definitions





\newcommand{\arctg}{\operatorname{arctg}}


\newcommand{\vraisup}{\operatorname{vrai~sup}}


\newcommand{\tts}[1]{}


\renewcommand{\baselinestretch}{1.1}





%\hoffset=-15mm                  For Eng.





\setcounter{page}{33}


\god{1997}


\nomer{11~(426)} %                Remove in Eng.


%\nomer{Vol.\,41, No.\,11}        Set in Eng.


% \str{75--81}                    Set in Eng.


\udk{517.988}


\author{\it  Т.М.\,Леденева}


\title{Некоторые аспекты представления нечетких операторов отношением двух многочленов}





%\thanks{Работа выполнена при поддержке РФФИ, грант \No. 94-01-00183.}


\date{}


% \pagestyle{myheadings}         Set in Eng.





\pagestyle{plain} %              Remove in Eng.


\begin{document}





% \thispagestyle{plain}          Set in Eng.





\maketitle





% Body text





Агрегирование нечеткой информации в системах принятия решений, которые


разрабатываются как составная часть экспертных систем, связано с


использованием нечетких операторов объединения и пересечения, поэтому


актуальными являются исследования, связанные с выявлением некоторых


свойств этих операторов, полезных для приложений, их структуры, а также


взаимосвязи различных представлений. Решение этих задач позволит


осуществить целенаправленный выбор оператора и тем самым повысить уровень


адекватности соответствующей модели.





Среди нечетких операторов, упоминаемых в литературе, можно выделить класс


операторов объединения и пересечения, которые представляются отношением


многочленов или, в частном случае, многочленом. К таким операторам


относятся следующие:


\begin{gather*}


F_m(x,y)=\max(x+y-1,0);\qquad G_m(x,y)=\min(x+y,1);\\


F_p(x,y)=xy;\qquad G_p(x,y)=x+y-xy;\\


F_0(x,y)=\frac{xy}{x+y-xy},\quad F_0(0,0)=0;\qquad


G_{-1}(x,y)=\frac{x+y-2xy}{1-xy},\quad G_{-1}(1,1)=1;\\


F_\alpha(x,y)=\frac{xy}{\alpha+(1-\alpha)(x+y-xy)};\qquad


G_\beta(x,y)=\frac{(\beta-1)xy+x+y}{1+\beta xy}.


\end{gather*}


Здесь $F_m$, $F_p$, $F_0$, $F_\alpha$ --- операторы пересечения,


$G_m$, $G_p$,


$G_{-1}$, $G_\beta$ --- операторы объединения (\cite{Tri}, с.221).





Нечеткие операторы обобщаются представлением их в классе треугольных норм


и конорм, причем операции пересечения определяются через $T$-нормы, а


объединения --- через $S$-конормы (\cite{Pos}, с.30).





Треугольная $T$-норма есть операция $T:[0,1]\times[0,1]\to[0,1]$,


удовлетворяющая условиям


коммутативности, ассоциативности, монотонности и ограниченности


$T(0,0)=0$, $T(1,x)=x$. Ее


конорма ($S$-конорма) определяется соотношением


$S(x,y)=1-T(1-x,1-y)$, удовлетворяет тем же


условиям, но ограниченность задается в виде


$S(1,1)=1$, $S(0,x)=x$.





Пусть


\begin{equation}


\widetilde\Phi(x,y)=\frac{a_0+a_1x+a_2y+a_3xy}


{b_0+b_1x+b_2y+b_3xy}


\end{equation}


реализует операцию объединения или пересечения нечетких множеств.





Рассмотрим свойства коммутативности и ассоциативности для


$\widetilde\Phi(x,y)$, которые


запишутся соответственно в виде:


\begin{align}


\widetilde\Phi(x,y)&=\widetilde\Phi(y,x),\\


\widetilde\Phi(\widetilde\Phi(x,y),z)&=


\widetilde\Phi(x,\widetilde\Phi(y,z)).


\end{align}





Найдем условия на коэффициенты


$\widetilde\Phi(x,y)$, при которых выполняются (2) и (3).


Подставляя $\widetilde\Phi(x,y)$


в виде (1) в (2) и (3), а затем приравнивая после


соответствующих преобразований коэффициенты при одинаковых степенях


$x$ в


левой и правой частях выражений (2) и (3), получим ряд соотношений,


анализ которых позволяет сделать следующие выводы:





$\widetilde\Phi(x,y)$ коммутативна при выполнении соотношений


\begin{align}


&


\begin{cases}


a_1=a_2,\\


b_1=b_2


\end{cases}


\\


\intertext{и ассоциативна при выполнении соотношений}


&


\begin{cases}


a_1=a_2,\quad b_1=b_2;\\


b_1a_0+a^2_1=a_1b_0+a_3b_0;\\


b_3a_1+b^2_1=b_0b_3+a_3b_1;\\


a_1b_1=a_0b_3.


\end{cases}


\end{align}





Учитывая соотношения (4), которые также входят в (5), в дальнейшем


коммутативную и ассоциативную операцию будем рассматривать в виде


\begin{equation}


\Phi(x,y)=\frac{a_0+a_1(x+y)+a_2xy}{b_0+b_1(x+y)+b_2xy},


\end{equation}


а соотношения (5) при этом запишутся следующим образом:


$$


\begin{cases}


b_1a_0+a^2_1=a_1b_0+a_2b_0,\\


b_2a_1+b^2_1=b_0b_2+a_2b_1,\\


a_1b_1=a_0b_2.


\end{cases}


$$





Известно, что ассоциативная операция представима в виде (\cite{Pos},


с.30)


$$


\Phi(x,y)=\varphi^{-1}(\varphi(x)+\varphi(y)),


$$


где $\varphi:[0,1]\to R^+$ --- непрерывная монотонная функция,


определяемая с точностью


до положительной константы и называемая аддитивным генератором;


$\varphi^{-1}$ ---


псевдообратная функция такая, что


$\varphi^{-1}(\varphi(x))=x$.





Согласно Де Финетти, если $\Phi(x,y)$ симметрична (обладает свойством


коммутативности) и имеет производную, то необходимым и достаточным


условием для ее ассоциативности является то, чтобы отношение двух частных


производных $\Phi$ было отношением функции одного $x$ к той же функции


одного $y$. В этом случае $\varphi$ определяется следующим образом:


\begin{equation}


\frac{\Phi'_x(x,y)}{\Phi'_y(x,y)}=


\frac{\varphi'(x)}{\varphi'(y)}.


\end{equation}





Рассмотрим задачу: {\it как для известной $\Phi(x,y)$ в виде {\em (6)}


определить аддитивный генератор $\varphi(x)$}?





Заметим, что для $\Phi(x,y)$ в виде (6) справедливо (7) с


$$


\varphi'(x)=\frac{1}


{(a_1b_0-a_0b_1)+x(a_2b_0-a_0b_2)+x^2(a_2b_1-a_1b_2)},


$$


тогда


\begin{equation}


\varphi(x)=\int\frac{dx}


{(a_1b_0-a_0b_1)+x(a_2b_0-a_0b_2)+x^2(a_2b_1-a_1b_2)}+C.


\end{equation}





В зависимости от коэффициентов $\Phi(x,y)$ возможны следующие ситуации


(табл.~1),


в каждой из которых аддитивный генератор получается в результате


вычисления соответствующего неопределенного интеграла.





{\hspace{107mm} Таблица 1}





{\hfill


\begin{tabular}{|c|c|c|c|c|c|c|c|c|}


\hline


Ситуация & $S_1$ & $S_2$ & $S_3$ & $S_4$ & $S_5$ & $S_6$


& $S_7$ & $S_8$\\


\hline


$\alpha=a_1b_0-a_0b_1$ &$+$ &$+$ &$+$ &$+$ &$-$ &$-$ &$-$ &$-$\\


\hline


$\beta=a_2b_0-a_0b_2$ &$+$ &$+$ &$-$ &$-$ &$+$ &$+$ &$-$ &$-$\\


\hline


$\gamma=a_2b_1-a_1b_2$ &$+$ &$-$ &$+$ &$-$ &$+$ &$-$ &$+$ &$-$\\


\hline


\end{tabular}


\hfill}


\medskip





\noindent (Здесь ``$+$'' означает, что выражение равно 0, ``$-$'' ---


иначе.)





Рассмотрим, например, ситуацию $S_4$. В этом случае


\begin{multline*}


\varphi(x)=\int\frac{dx}


{x(a_2b_0-a_0b_2)+x^2(a_2b_1-a_1b_2)}=


\int\frac{dx}{\beta x+\gamma x^2} =\frac{1}{\gamma}\int


\frac{dx}{x(x+\frac{\beta}{\gamma})}=\\


=\frac{1}{\beta}\biggl(\int\frac{dx}{x}-\int


\frac{dx}{x+\frac{\beta}{\gamma}}\biggr)=\frac{1}{\beta}


\Bigl(\ln|x|-\ln\bigl|x+\frac{\beta}{\gamma}\Bigr|\Bigr)=


\frac{1}{\beta}\ln\biggl|\frac{x}{x+\frac{\beta}{\gamma}}\biggr|=


\frac{1}{\beta}\ln\Bigl|\frac{\gamma x}{\gamma x+\beta}\Bigr|+C.


\end{multline*}





Таким образом, для каждой из перечисленных ситуаций имеем





$S_1$: аддитивного генератора не существует;





$S_2$: $\varphi(x)=-\dfrac{1}{\gamma x}+C$;





$S_3$: $\varphi(x)=\dfrac{1}{\beta}\ln|x|+C$;





$S_4$: $\varphi(x)=\dfrac{1}{\beta}\ln\Bigl|


\dfrac{\gamma x}{\gamma x+\beta}\Bigr|+C$;





$S_5$: $\varphi(x)=\dfrac{1}{\alpha}x+C$;





$S_6$: если $\dfrac{\alpha}{\gamma}>0$, то


$\varphi(x)=\dfrac{1}{\sqrt{\alpha\gamma}}


\arctg\Bigl(x\sqrt{\dfrac{\gamma}{\alpha}}\Bigr)+C$, иначе


$\varphi(x)=\dfrac{1}{\sqrt{-\frac{\alpha}{\gamma}}}\ln


\Biggl|\dfrac{x-\sqrt{-\frac{\alpha}{\gamma}}}


{x+\sqrt{-\frac{\alpha}{\gamma}}}\Biggr|+C$;





$S_7$: $\varphi(x)=\dfrac{1}{\beta}\ln|\alpha+\beta x|+C$;





$S_8$: пусть $D=\beta^2-4\alpha\gamma$, тогда возможны случаи


\begin{itemize}


\item[\ ] если $D>0$, то


$\varphi(x)=\dfrac{1}{\sqrt D}\ln\Bigl|


\dfrac{2\gamma x+\beta-\sqrt D}


{2\gamma x+\beta+\sqrt D}\Bigr|+C$,\\


\item[\ ] если $D=0$, то


$\varphi(x)=-\dfrac{1}{\gamma x+\frac{\beta}{2}}+C$,\\


\item[\ ] если $d<0$, то


$\varphi(x)=\dfrac{2}{\sqrt{-D}}\arctg


\dfrac{2\gamma x+\beta}{\sqrt{-D}}+C$.


\end{itemize}





При использовании приведенных формул следует иметь в виду, что


$x\in[0,1]$, а


константа $C$ находится из дополнительных условий: в случае операции


пересечения аддитивный генератор должен быть монотонно убывающей


функцией, для которой $\varphi(1)=0$; для операции объединения генератор


представляется монотонно возрастающей функцией, причем


$\varphi(0)=0$.





В качестве примера рассмотрим операцию объединения


$$


G_{-1}(x,y)=\frac{x+y-2xy}{1-xy}.


$$


Здесь $a_0=0$, $a_1=1$, $a_2=-2$, $b_0=1$, $b_1=0$, $b_2=-1$,


$\alpha=1$, $\beta=-2$, $\gamma=1$.


Следовательно, имеем ситуацию $S_8$ с $D=0$. Тогда


$$


\varphi(x)=-\frac{1}{\gamma x+\dfrac{\beta}{2}}+C=


\frac{1}{1-x}+C.


$$


Из условия $\varphi(0)=0$ определяется константа $C=-1$, и аддитивный


генератор имеет вид


$$


\varphi(x)=\frac{x}{1-x}=g_{-1}(x),


$$


что соответствует результату, изложенному в (\cite{Tri}, с.221).





Возможность выписать в явном виде аддитивный генератор позволяет решить


ряд проблем, среди них --- определение операции отрицания, с помощью


которой для нечетких операторов объединения и пересечения вводится


понятие двойственности -- $n$-дуальности (\cite{Tri}, с.222). Операция


отрицания $n(x)$ на $[0,1]$, которая в общем случае может быть не


единственной, определяется следующим образом:


$$


n=f^{-1}\circ l_t\circ g=g^{-1}\circ l_{1/t}\circ f,


$$


где $g$, $f$ --- генераторы объединения и пересечения соответственно,


$l_t(x)=tx$


($t>0$).





Рассмотрим известные пары двойственных операторов (\cite{Tri}, с.222) и


выпишем операции отрицания. Для $F_p(x,y)$, $G_p(x,y)$ генераторы


$g_p(x)=-\ln(1-x)$ и $f_p(x)=-\ln x$


определяют


единственную функцию сильного отрицания $n(x)=1-x$. Для


$F_0(x,y)$, $G_{-1}(x,y)$ генераторы


$f_0(x)=\dfrac{1-x}{x}$ и $g_{-1}(x)=\dfrac{x}{1-x}$


определяют множество функций сильного отрицания Суджено при


$t>0$ в виде $n(x)=\dfrac{1-x}{1+(t-1)x}$.


Для $F_\alpha(x,y)$ и $G_\beta(x,y)$


имеем единственное отрицание Суджено


$$


n(x)=\frac{1-x}{1+\dfrac{1-\alpha+\beta}{\alpha}x},


$$


полученное с помощью генераторов


$f_\alpha(x)=\ln\dfrac{\alpha-(\alpha-1)x}{x}$ и


$g_\beta(x)=\ln\dfrac{1+\beta x}{1-x}$.





Как видно из полученных результатов, аддитивный генератор для


$\Phi(x,y)$ в виде


(6) относится к одному из следующих типов:


$$


\varphi_1(x)=\frac{ax+b}{cx+d},\quad


\varphi_2(x)=\ln\frac{ax+b}{cx+d},\quad


\varphi_3(x)=\arctg\frac{ax+b}{cx+d}.


$$





Рассмотрим обратную задачу: {\it как с помощью известного генератора


$\varphi(x)$


построить $\Phi(x,y)$ в виде {\em (6)}}?





При решении этой задачи получены выражения для коэффициентов


$a_0$, $a_1$, $a_2$,


$b_0$, $b_1$, $b_2$ через коэффициенты аддитивных генераторов, причем





$\varphi_1(x)=\dfrac{ax+b}{cx+d}$ определяет


коммутативную и ассоциативную операцию $\Phi(x,y)$ с коэффициентами


$a_0=bd^2$,


$a_1=ad^2$, $a_2=2adc-bc^2$,


$b_0=ad^2-2bcd$, $b_1=-bc^2$, $b_3=-ac^2$;





$\varphi_2(x)=\ln\dfrac{ax+b}{cx+d}$ генерирует коммутативную и


ассоциативную операцию $\Phi(x,y)$ и, если $ad\ne bc$, то


\begin{alignat*}{3}


& a_0=\frac{db^2-bd^2}{ad-bc}, &\quad


& a_1=\frac{bd(a-c)}{ad-bc}, &\quad


& a_2=\frac{a^2d-bc^2}{ad-bc},\\


& b_0=\frac{ad^2-cb^2}{ad-bc}, &\quad


& b_1=\frac{ac(d-b)}{ad-bc},& \quad


& b_2=\frac{ac^2-a^2c}{ad-bc};


\end{alignat*}





$\varphi_3(x)=\arctg\dfrac{ax+b}{cx+d}$ определяет коэффициенты $\Phi(x,y)$


следующим образом:


\begin{alignat*}{3}


& a_0=b(d^2+b^2), &\quad & a_1=a(d^2+b^2), &\quad


& a_2=2acd-b(c^2-a^2),\\


& b_0=a(d^2-b^2)-2bcd, &\quad & b_1=-b(a^2+c^2), &\quad


& b_2=-a(a^2+c^2).


\end{alignat*}





Практический смысл результатов, полученных при решении сформулированных


выше задач, заключается в том, что, используя полученные соотношения


между коэффициентами $\Phi(x,y)$ и $\varphi(x)$, для любой коммутативной


и ассоциативной


операции, представленной выражением (6), можно определить аддитивный


генератор, а также для любого аддитивного генератора найти


соответствующую ему операцию $\Phi(x,y)$ в виде (6).





Определим теперь условия, при которых $\Phi(x,y)$ задается как


$T$-норма,


т.е. помимо коммутативности, ассоциативности и монотонности выполняются


соотношения


\begin{equation}


\begin{cases}


\Phi(0,0)=0,\\


\Phi(1,x)=x.


\end{cases}


\end{equation}





Рассмотрим различные типы аддитивных генераторов.





Используя представление коэффициентов $\Phi(x,y)$ через коэффициенты


$a$, $b$, $c$, $d$ генератора $\varphi_1(x)$, получим, что при


$c\ne 0\;$ (9)


справедлива при


$


\begin{cases}


d\ne 0,\\


a+b=0,


\end{cases}


$


что обусловливает единственность генератора $f_0(x)$ и соответствующей


ему


операции пересечения нечетких множеств $F_0(x,y)$, доопределенной


соотношением $F_0(0,0)=0$.





При $c=0\;$ $\varphi_1(x)$ представляется линейной функцией, и для


выполнения условий (9)


необходимо, чтобы $a+b=0$, что соответствует


$f_m(x)=1-x$ и оператору $\Phi(x,y)=x+y-1$, проектируя


который на $[0,1]$ получим $F_m(x,y)$.





В случае $\varphi_2(x)$ соотношения (9) справедливы при


$


\begin{cases}


a=d=0,\\


b=c\ne 0


\end{cases}


$


с аддитивным генератором $f_p(x)$ и операцией пересечения $F_p(x,y)$,


а также для


$


\begin{cases}


d=0,\\


a+b=c,\\


bc\ne 0,


\end{cases}


$


при этом $\varphi(x)=\ln\dfrac{ax+b}{(a+b)x}$


$(a/b>-1)$ порождает пересечение вида


$\Phi(x,y)=\dfrac{(a+b)xy}{b+a(x+y-xy)}$. Положив


$\dfrac{b}{a+b}=\alpha$, получим $f_\alpha(x)$ и $F_\alpha(x,y)$.





Определим теперь условия, при которых $\Phi(x,y)$ задается как


$S$-конорма, т.е.


наряду с условиями коммутативности, ассоциативности и монотонности,


выполняются следующие:


\begin{equation}


\begin{cases}


\Phi(1,1)=1,\\


\Phi(0,x)=x.


\end{cases}


\end{equation}





В случае $\varphi_1(x)$ соотношения (10) выполняются при


$


\begin{cases}


b=0,\\


c+d=0,


\end{cases}


$


что соответствует аддитивному генератору


$g_{-1}(x)$ и операции объединения $G_{-1}(x,y)$,


доопределенной соотношением $G_{-1}(1,1)=1$.





Рассматривая $\varphi_1(x)$ при $c=0$, можно показать, что при


$ad\ne 0$ условия


(10)


справедливы только при $b=0$, тогда имеем $g_m(x)$, порождающий


$\Phi(x,y)=x+y$,


проектируя который на $[0,1]$, получим $G_m(x,y)$.





Для $\varphi_2(x)$ существует несколько возможностей выполнения (10).


\begin{itemize}


\item[i)]


$


\begin{cases}


a=0\\


d=b=-c


\end{cases}


$


определяет аддитивный генератор $g_p(x)$ и соответствующую ему операцию


объединения $G_p(x,y)$;\\


\item[ii)]


$


\begin{cases}


a\ne c\\


d=b=-c


\end{cases}


$


соответствует генератору


$\varphi(x)=\ln\dfrac{ax+b}{b(1-x)}$ $(a/b>1)$


и операции объединения


$\Phi(x,y)=\dfrac{b(x+y)+(a-b)xy}{b+axy}$. Если положить


$a/b=\beta$,


то получим $g_\beta(x)$ и $G_\beta(x,y)$;\\


\item[iii)]


$b=d=a=-c$ позволяет получить аддитивный генератор


$\varphi(x)=\ln\dfrac{1+x}{1-x}$ и операцию


объединения $G_l(x,y)=\dfrac{x+y}{1+xy}$, известную под названием


оператора Лоренца.


\end{itemize}





$\varphi_3(x)=\arctg\dfrac{ax+b}{cx+d}$ ни при каких условиях на


коэффициенты $a$, $b$, $c$, $d$


не порождает треугольные нормы и конормы.





Анализ условий, накладываемых на коэффициенты аддитивных генераторов,


позволяет сделать вывод о невозможности получить другие представления


нечетких операторов объединения и пересечения, отличные от известных (в


том числе параметрических) в виде отношения двух многочленов (в частном


случае --- многочленом).





Исследуем взаимосвязи, существующие между различными представлениями


нечетких объединений и пересечений. Выберем в качестве базовых операций


на $[0,1]$ алгебраическую сумму


$x\oplus y=x+y-xy$ ($x\oplus 0=x$, $x\oplus 1=1$) и алгебраическое


произведение


$xy$ ($x\cdot 0=0$, $x\cdot 1=1$) и покажем, что с их помощью могут быть


выражены рассматриваемые нечеткие операторы.





Заметим, что $G_p(x,y)=x\oplus y$, $F_p(x,y)=xy$. 





Если


$x,y\in[0,1]$, $x+y-xy=z$, то


$$


x=\frac{z-y}{1-y}=z\theta y,\qquad


y=\frac{z-x}{1-x}=z\theta x,


$$


где $\theta$ --- операция вычитания для операции сложения $\oplus$,


которую будем


называть нормированной разностью между $x$ и $y$. Справедливы следующие


свойства операции $\theta$:


$$


x\theta 0=x;\quad x\theta x=0;\quad 1\theta x=1;\quad


x\theta y\ne y\theta x.


$$





Заметим, что операция вычитания $\theta$ связана с процессом


нормирования,


представляющим собой линейное преобразование, при котором область


значений любого ограниченного действительного отображения $f$


переводится в интервал $[0,1]$, т.е.


$$


f\to (f-\inf f)/(\sup f-\inf f).


$$





Пусть


$t(x_{\min},x,x_{\max})=\dfrac{x-x_{\min}}{x_{\max}-x_{\min}}$


--- функция нормирования, заданная на $[x_{\min},x_{\max}]$, тогда на


$[0,1]$ при


$0<x<y<1\;$ $t(0,x,y)=\dfrac{x}{y}$ --- это нормированное значение $x$ на


$[0,y]$; $t(x,y,1)=\dfrac{x-y}{1-y}=x\theta y$ --- это нормированное


значение $x$ на $[y,1]$. \;$[0,y]$ и $[y,1]$ будем называть интервалами


нормирования.





Отметим свойства функции $t(x_{\min},x,x_{\max})$.





{\it Свойство} 1. Если $0<x<x_1<x_2$, то


$t(0,x,x_2)=t(0,x,x_1)\cdot t(0,x_1,x_2)$, т.е. при расширении интервала


нормирования


новое нормированное значение представляет собой произведение старого


нормированного значения и нормированного значения прежней границы на


новом интервале нормирования.





{\it Свойство} 2. Если $x_2<x_1<x<1$, то


$t(x_1,x,1)\oplus t(x_2,x_1,1)=t(x_2,x_1,1)$ ---


при расширении интервала


нормирования, т.е.


если вместо $x_1$ выбрать значение $x_2<x_1$, то нормированное значение


$x$ на $[x_2,1]$


представляет собой алгебраическую сумму нормированного значения


$x$ на $[x_1,1]$


и нормированного значения прежней границы $x_1$ на новом интервале


нормирования $[x_2,1]$.





Пусть $x,y\in[0,1]$, тогда $xy\le x\oplus y$. Можно показать, что


\begin{align*}


F_0(x,y)&=\frac{xy}{x+y-xy}=\frac{xy}{x\oplus y}=


t(0,xy,x\oplus y);\\


G_{-1}(x,y)&=\frac{x+y-2xy}{1-xy}=


\frac{(x\oplus y)-xy}{1-xy}=(x\oplus y)\theta xy


=t(xy,x\oplus y,1);\\


F_\alpha(x,y)&=


\frac{xy}{\alpha+(1-\alpha)(x+y-xy)}=


\frac{xy}{\alpha\oplus(x\oplus y)}=


t(0,xy,\alpha\oplus(x\oplus y))\quad (\alpha>0);\\


G_\beta(x,y)&=


\frac{(\beta-1)xy+x+y}{1+\beta xy}=


\frac{(x\oplus y)+\beta xy}{1+\beta xy}=


(x\oplus y)\theta(-\beta xy)=


t(-\beta xy,x\oplus y,1)\quad (\beta>-1).


\end{align*}





Таким образом, $F_0$ и $F_\alpha$, моделирующие пересечение нечетких


множеств,


определяют нормированное значение $xy$ на $[0,x\oplus y]$ или на


интервале нормирования


с расширенной правой границей $[0,\alpha\oplus(x\oplus y)]$.


\;$G_{-1}$ и $G_\beta$,


моделирующие объединение


нечетких множеств, определяют нормированное значение


$x\oplus y$ на интервале


нормирования $[xy,1]$ или на интервале нормирования с расширенной левой


границей $[-\beta xy,1]$.





При $\alpha=0\;$ $F_\alpha(x,y)=F_0(x,y)$, при


$\beta=-1\;$ $G_\beta(x,y)=G_{-1}(x,y)$.





В соответствии со свойством 1


$$


F_\alpha(x,y)=t(0,xy,\alpha\oplus(x\oplus y))=


t(0,xy,x\oplus y)t(0,x\oplus y,\alpha\oplus(x\oplus y))=


F_0(x,y)\frac{x\oplus y}{\alpha\oplus(x\oplus y)}.


$$





Заметим, что при $\alpha>0$ для фиксированных значений $x,y\in[0,1]$


$$


F_\alpha(x,y)<F_0(x,y)<F_p(x,y),


$$


причем при $\alpha\to\infty\;$ $F_\alpha(x,y)\to 0$.





В соответствии со свойством 2


$$


G_\beta(x,y)=t(-\beta xy,x\oplus y,1)=t(xy,x\oplus y,1)\oplus


t(-\beta xy,xy,1)=G_{-1}(x,y)\oplus


\frac{xy+\beta xy}{1+\beta xy}.


$$





Заметим, что при $\beta>-1$ для фиксированных значений $x,y\in[0,1]$


$$


G_\beta(x,y)>G_{-1}(x,y)>G_p(x,y),


$$


причем при $\beta\to\infty\;$ $G_\beta(x,y)\to 1$.





Учитывая, что $F_p(x,y)<G_p(x,y)$, получим следующие соотношения между


нечеткими операторами объединения и пересечения:


$$


F_\alpha(x,y)<F_0(x,y)<F_p(x,y)<G_p(x,y)<


G_{-1}(x,y)<G_\beta(x,y).


$$





Важное значение для выбора вида оператора в прикладных задачах имеет


характеристика носителя и тип задачи. Если рассматривается выбор


наилучшего варианта на насыщенном носителе (избыток претендентов), то


целесообразно расширить интервал нормирования, что обусловливает


использование операторов объединения и пересечения с высоким уровнем


жесткости (увеличение $\alpha$ для $F_\alpha$ и уменьшение


$-\beta$ для $G_\beta$). На


обедненном


носителе уровень жесткости следует снижать. В задаче ранжирования


вариантов наряду с $F_p$ и $G_p$ возможно использование


$F_\alpha$ и $G_\beta$ с


небольшим


уровнем жесткости. В особенно сложных задачах принятия решений необходимо


провести предварительное исследование, и только на основе тщательного


анализа задачи и свойств операторов следует выбрать тот, который наиболее


адекватно отражает взаимодействие нечетких множеств.








\begin{thebibliography}{9}





\bibitem{Tri}


Трильяс Э., Альсина К., Вальверде Л. {\em Нужны ли в теории нечетких


множеств операции $\max$, $\min$, $1-j$}? // Нечеткие множества и теория


возможностей. Последние достижения: Пер. с англ. / Под ред.


Р.Р.\,Ягера. -- М.: Радио и связь, 1986. -- 405 с.


\bibitem{Pos}


{\em Нечеткие множества в моделях управления и искусственного


интеллекта} / Под ред. Поспелова Д.А. -- М.: Наука, 1986. -- 312 с.








\end{thebibliography}





\vskip 2cm





\postup{\it Воронежский государственный университет}{$\qquad$\it


Поступили}


\postup{}{{\it первый вариант} 30.01.1995}


\postup{}{{\it окончательный вариант} 20.05.1997}





\end{document}


